\chapter{Simulated Quantum Annealing (\code{sqa})}
\label{ch:sqa}

This chapter contains the theoretical heart of \qv. We introduce the
transverse-field Ising model, state the adiabatic theorem that motivates
quantum annealing, and then derive --- in full detail --- the Suzuki--Trotter
path-integral mapping that turns the quantum partition function into a
classical Ising model in one extra dimension. The result is the exact
effective action and acceptance rules executed by the C++
\code{SQAAnnealer}, including the Trotter coupling
$\Jp=\tfrac12\ln\coth(\beta\Gamma/M)$ that reappears throughout the rest of
the manual.

\section{Quantum annealing in one page}
\label{sec:qa}

Quantum annealing~\cite{kadowaki,farhi} encodes the optimisation problem in
the \emph{longitudinal} part of a quantum spin Hamiltonian and adds a
\emph{transverse field} that induces quantum fluctuations:

\begin{keyeq}[Transverse-field Ising model (TFIM)]
\begin{equation}
\hat H(\Gamma)
=\underbrace{\sum_i h_i\,\hat\sigma^z_i+\sum_{i<j}J_{ij}\,\hat\sigma^z_i\hat\sigma^z_j}_{\hat H_{\mathrm{cl}}
\;\text{(the problem)}}
\;-\;\Gamma\sum_i \hat\sigma^x_i ,
\label{eq:tfim}
\end{equation}
\end{keyeq}

\noindent where $\hat\sigma^{x,z}_i$ are Pauli operators on spin $i$.
$\hat H_{\mathrm{cl}}$ is diagonal in the computational ($\sigma^z$) basis
and its ground state is exactly the Ising optimum we seek. The transverse
term $-\Gamma\sum_i\hat\sigma^x_i$ flips spins
($\hat\sigma^x\ket{\pm1}=\ket{\mp1}$) and therefore lets the state
\emph{tunnel} between classical configurations.

The protocol: start at large $\Gamma$, where the ground state is the trivial
uniform superposition $\bigotimes_i\tfrac{1}{\sqrt2}(\ket{+1}+\ket{-1})$;
reduce $\Gamma\to0$ slowly. The \textbf{adiabatic theorem} guarantees that if
the sweep is slow compared to the inverse square of the spectral gap
$\Delta_{\mathrm{gap}}(\Gamma)$ (the energy difference between the ground and
first excited state), the system remains in its instantaneous ground state
throughout and ends in the classical optimum. The total time required scales
as
\begin{equation}
T \;\gtrsim\; \frac{\max_\Gamma\bigl|\bra{1}\partial_\Gamma\hat H\ket{0}\bigr|}{\min_\Gamma \Delta_{\mathrm{gap}}^2(\Gamma)} ,
\label{eq:adiabatic}
\end{equation}
so the bottleneck is the \emph{minimum gap}, typically encountered at a
quantum phase transition separating the disordered (large $\Gamma$) phase
from the ordered (small $\Gamma$) phase. This observation drives the optimal
schedule of Chapter~\ref{ch:optimal}: sweep slowly \emph{only} near the
transition.

\begin{physbox}
Thermal hopping surmounts a barrier of height $B$ at rate
$\sim\mathrm{e}^{-\beta B}$, indifferent to the barrier's width. Quantum
tunnelling through a barrier of width $w$ and height $B$ proceeds at rate
$\sim\mathrm{e}^{-c\,w\sqrt{B}\,/\Gamma}$: it pays for \emph{width}, not
height. Rugged landscapes whose minima are separated by tall-but-thin walls
are therefore the natural habitat of quantum annealing --- and of its
classical-hardware simulation, SQA.
\end{physbox}

\section{The Suzuki--Trotter mapping, in full}
\label{sec:trotter}

Real-time simulation of Eq.~\eqref{eq:tfim} on a CPU is exponentially
expensive. \emph{Simulated} quantum annealing instead samples the quantum
system's \emph{thermal equilibrium} at inverse temperature $\beta$ ---
the density matrix $\mathrm{e}^{-\beta\hat H}/Z$ --- using a mapping to a
classical model that Monte Carlo can handle. We now derive it.

\subsection{Step 1: split the exponential}

The partition function is $Z=\Tr\,\mathrm{e}^{-\beta\hat H}$ with
$\hat H=\hat H_{\mathrm{cl}}+\hat H_q$,
$\hat H_q=-\Gamma\sum_i\hat\sigma^x_i$. Since
$[\hat H_{\mathrm{cl}},\hat H_q]\ne0$ we cannot factorise the exponential
exactly, but we can \emph{Trotterise}: write
$\mathrm{e}^{-\beta\hat H}=\bigl(\mathrm{e}^{-\frac{\beta}{M}\hat H}\bigr)^M$
and apply the Lie--Trotter splitting to each thin slice,
\begin{equation}
\mathrm{e}^{-\frac{\beta}{M}(\hat H_{\mathrm{cl}}+\hat H_q)}
=\mathrm{e}^{-\frac{\beta}{M}\hat H_{\mathrm{cl}}}\,
 \mathrm{e}^{-\frac{\beta}{M}\hat H_q}
 +O\!\left(\frac{\beta^2}{M^2}\bigl\|[\hat H_{\mathrm{cl}},\hat H_q]\bigr\|\right).
\label{eq:trotter-split}
\end{equation}
The neglected term is the \emph{Trotter error}; it vanishes as
$M\to\infty$ like $O((\beta/M)^2)$ per slice, i.e.\ $O(\beta^2/M)$ overall.
$M$ is the \code{trotter\_slices} parameter.

\subsection{Step 2: insert complete sets of classical states}

Insert the identity
$\mathbb 1=\sum_{s}\ket{s}\bra{s}$ (sum over all $2^n$ classical spin
configurations) between every pair of factors. Writing $s^t$ for the
configuration inserted at slot $t=1,\dots,M$ and using cyclicity of the
trace ($s^{M+1}\equiv s^1$, \emph{periodic boundary in imaginary time}):
\begin{equation}
Z\;\approx\;\sum_{s^1,\dots,s^M}\;\prod_{t=1}^{M}
\underbrace{\bra{s^t}\mathrm{e}^{-\frac{\beta}{M}\hat H_{\mathrm{cl}}}\ket{s^t}}_{\displaystyle
\mathrm{e}^{-\frac{\beta}{M}\Ecl(s^t)}}
\;\underbrace{\bra{s^t}\mathrm{e}^{\frac{\beta\Gamma}{M}\sum_i\hat\sigma^x_i}\ket{s^{t+1}}}_{\text{transverse matrix element}} .
\label{eq:z-inserted}
\end{equation}
The classical factor is diagonal, giving the Boltzmann weight of each slice
at the \emph{reduced} inverse temperature $\beta/M$. The transverse factor
factorises over sites, so we only need the $2\times2$ single-spin matrix
element.

\subsection{Step 3: the single-spin transfer matrix}

Let $a=\beta\Gamma/M$. Using
$\mathrm{e}^{a\hat\sigma^x}=\cosh a\,\mathbb 1+\sinh a\,\hat\sigma^x$:
\begin{equation}
\bra{s}\mathrm{e}^{a\hat\sigma^x}\ket{s'}
=\begin{cases}
\cosh a & s=s'\\
\sinh a & s\ne s' .
\end{cases}
\end{equation}
We now rewrite this as a classical Ising bond between the same site in
adjacent slices, i.e.\ find $\Lambda$ and $\Jp$ such that
$\bra{s}\mathrm{e}^{a\hat\sigma^x}\ket{s'}=\Lambda\,\mathrm{e}^{\Jp\,ss'}$.
Matching the two cases ($ss'=+1$ and $ss'=-1$):
\begin{equation}
\Lambda\,\mathrm{e}^{\Jp}=\cosh a,\qquad
\Lambda\,\mathrm{e}^{-\Jp}=\sinh a
\quad\Longrightarrow\quad
\mathrm{e}^{2\Jp}=\coth a,\qquad
\Lambda^2=\sinh a\cosh a=\tfrac12\sinh 2a .
\end{equation}

\begin{keyeq}[The Trotter coupling]
\begin{equation}
\Jp(\beta,\Gamma,M)\;=\;\frac12\,\ln\coth\!\Bigl(\frac{\beta\Gamma}{M}\Bigr)
\;=\;\frac12\,\ln\!\frac{1}{\tanh(\beta\Gamma/M)}\;>\;0 .
\label{eq:jperp}
\end{equation}
\end{keyeq}

$\Jp$ is a \emph{ferromagnetic} bond along imaginary time (it rewards
$s^t_i=s^{t+1}_i$; note the effective action below carries it with a minus
sign). Its two limits define the whole phenomenology:
\begin{itemize}
\item $\Gamma\to\infty$ (strong quantum fluctuations): $a$ large,
$\coth a\to1$, $\Jp\to0$ --- slices decouple, worldlines flap freely.
\item $\Gamma\to0^+$ (classical limit): $a\to0$, $\Jp\to\infty$ --- slices
lock together into identical copies of one classical configuration.
\end{itemize}
The Python helper \code{j\_perp\_from\_beta\_gamma(beta, gamma,
trotter\_slices)} evaluates Eq.~\eqref{eq:jperp} directly.

\begin{warnbox}[Direction of the ramp]
Because $\Jp$ is monotonically \emph{decreasing} in $\Gamma$, an annealing
run that lowers $\Gamma$ corresponds to \emph{raising} $\Jp$. All
\qv{} adaptive-schedule APIs are parameterised by $\Jp$ and therefore ramp it
\emph{upward}, from \code{j\_perp\_start} (quantum end) to
\code{j\_perp\_end} (classical end). Keep this inversion in mind when reading
traces.
\end{warnbox}

\subsection{Step 4: the effective classical action}

Substituting both factors into Eq.~\eqref{eq:z-inserted}:

\begin{keyeq}[SQA effective action]
\begin{equation}
Z\;\approx\;C\!\!\sum_{\{s^t_i\}}\mathrm{e}^{-S[s]},
\qquad
S[s]\;=\;\frac{\beta}{M}\sum_{t=1}^{M}\Ecl(s^t)
\;-\;\Jp\sum_{t=1}^{M}\sum_{i=1}^{n}s^t_i\,s^{t+1}_i ,
\label{eq:action}
\end{equation}
with $C=\Lambda^{nM}$ an irrelevant constant and periodic slices
$s^{M+1}\equiv s^1$.
\end{keyeq}

This is a purely classical Ising model on $n\times M$ spins: $M$ copies
(``Trotter slices'') of the original problem, each internally coupled by the
rescaled problem couplings $\frac{\beta}{M}J_{ij},\frac{\beta}{M}h_i$, and
neighbouring slices tied site-by-site with strength $\Jp$. Quantum
statistics at inverse temperature $\beta$ became classical statistics of a
$(d{+}1)$-dimensional system at unit temperature. Metropolis sampling of
$\mathrm{e}^{-S}$ (Chapter~\ref{ch:statmech}, with $S$ in place of
$\beta E$) is exact for this effective model.

\begin{figure}[ht]
\centering
\begin{tikzpicture}[scale=0.92,
  spin/.style={circle,draw=qnavy,fill=qnavy!12,inner sep=1.6pt,minimum size=8pt},
  intra/.style={qblue!75,thick},
  inter/.style={qteal,very thick,densely dashed}]
% slices: 4 slices x 5 spins
\foreach \t in {0,1,2,3}{
  \begin{scope}[shift={(\t*3.1,0)}]
    \foreach \i in {0,...,4}{ \node[spin] (s\t\i) at (0,\i*0.85) {}; }
    % intra-slice couplings (a few representative)
    \draw[intra] (s\t0) to[bend right=35] (s\t2);
    \draw[intra] (s\t1) to[bend left=35] (s\t4);
    \draw[intra] (s\t2) -- (s\t3);
    \draw[intra] (s\t0) -- (s\t1);
    \node[below=2mm of s\t0,font=\scriptsize\sffamily\color{qgray}]
      {slice $t{=}\pgfmathparse{int(\t+1)}\pgfmathresult$};
  \end{scope}}
% inter-slice bonds
\foreach \t/\tn in {0/1,1/2,2/3}{
  \foreach \i in {0,...,4}{ \draw[inter] (s\t\i) -- (s\tn\i); }}
% periodic wrap
\foreach \i in {0,...,4}{
  \draw[inter,opacity=0.45] (s3\i) .. controls +(1.2,0.55) and +(-1.2,0.55) .. (s0\i);}
% labels
\node[font=\small\sffamily\color{qblue}] at (1.55,4.6)
  {$\dfrac{\beta}{M}J_{ij}$ \; (problem couplings, within slice)};
\node[font=\small\sffamily\color{qteal}] at (7.75,-1.35)
  {$-\Jp\, s^t_i s^{t+1}_i$ \; (Trotter bonds along imaginary time, periodic)};
\end{tikzpicture}
\caption{The Suzuki--Trotter mapping: $M$ replicas of the classical problem
(here $M{=}4$, $n{=}5$) coupled site-by-site along a periodic imaginary-time
ring. The row of same-site spins across all slices is that site's
\emph{worldline}.}
\label{fig:trotter}
\end{figure}

\subsection{Trotter error and choosing $M$}

The systematic error of the discretisation is $O\!\left((\beta\Gamma/M)^2\right)$
per slice. In practice:
\begin{itemize}
\item $M=16$--$32$ (the default is 32) suffices for optimisation use, where
we care about ground-state \emph{identification} rather than unbiased
thermodynamics.
\item Increase $M$ if the per-slice energies disagree strongly with each
other late in the anneal (a Trotter-error signature), or when running at
large $\beta\Gamma$.
\item Memory and time grow linearly in $M$; the OpenMP path parallelises
over replicas and slices, softening the wall-clock cost.
\item \code{continuous\_time\_slices > 0} overrides $M$ with a much larger
value to approximate the continuous-time limit within the discrete engine;
the genuinely continuous algorithm is Chapter~\ref{ch:ctpimc}.
\end{itemize}

\section{Monte Carlo moves on the Trotter system}
\label{sec:sqa-moves}

The \code{SQAAnnealer} mixes three move types, each preserving detailed
balance for the action Eq.~\eqref{eq:action}. Every acceptance test is
$\min(1,\mathrm{e}^{-\Delta S})$; what differs is the update and its
$\Delta S$.

\subsection{Single-spin (slice) moves}

Flip one spin $s^t_i$. The action change has a classical part within the
slice and a quantum part from the two adjacent time-bonds:
\begin{keyeq}[Single-spin action change]
\begin{equation}
\Delta S=\frac{\beta}{M}\,\Delta\Ecl
\;+\;2\,\Jp\,s^t_i\bigl(s^{t-1}_i+s^{t+1}_i\bigr),
\qquad \Delta\Ecl=-2\,s^t_i f^t_i ,
\label{eq:sqa-flip}
\end{equation}
\end{keyeq}
where $f^t_i$ is the cached local field of slice $t$
(Eq.~\eqref{eq:deltaE}). Both pieces are $O(1)$; the implementation keeps
one local-field cache per (replica, slice). \code{sweeps\_per\_beta} counts
these sweeps.

\subsection{Worldline moves}

Flip site $i$ in \emph{every} slice simultaneously
($s^t_i\to-s^t_i\ \forall t$). All imaginary-time bonds
$s^t_is^{t+1}_i$ are products of two flipped spins and are therefore
\emph{invariant}: the Trotter term cancels exactly, leaving
\begin{equation}
\Delta S_{\text{worldline}}=\frac{\beta}{M}\sum_{t=1}^{M}\Delta\Ecl^{(t)} .
\label{eq:worldline}
\end{equation}
Worldline moves are the quantum analogue of a classical single-spin flip and
are essential late in the anneal: when $\Jp$ is large the slices are locked
and single-slice flips are strongly suppressed
($\Delta S\approx4\Jp$), while a worldline flip pays no Trotter penalty at
all. \code{worldline\_sweeps} (default 5 in \code{solve()}) counts these.

\subsection{Swendsen--Wang cluster moves along imaginary time}
\label{sec:sw-moves}

Between the two extremes sits the cluster move: flip a \emph{contiguous
segment} of a worldline. Segments are grown stochastically with the
Fortuin--Kasteleyn rule specialised to the 1-D ferromagnetic time-chain of
coupling $\Jp$: starting from a random seed $(i,t)$, extend the cluster
across an aligned time-bond with the \emph{join probability}

\begin{keyeq}[Cluster join probability]
\begin{equation}
p_{\text{join}}=1-\mathrm{e}^{-2\Jp}.
\label{eq:joinprob}
\end{equation}
\end{keyeq}

Bonds inside the cluster then flip for free (the FK construction makes the
grown-bond factor cancel from the acceptance ratio); the remaining action
change --- the classical part in every touched slice plus the two
\emph{boundary} time-bonds where the cluster ends --- is Metropolis-tested:
\begin{equation}
\Delta S_{\text{cluster}}
=\frac{\beta}{M}\!\!\sum_{t\in\text{cluster}}\!\!\Delta\Ecl^{(t)}
\;+\;2\,\Jp\!\!\sum_{\text{boundary bonds}}\!\! s\,s' .
\end{equation}
Cluster moves interpolate smoothly between single-spin flips (clusters of
size 1 when $\Jp$ is small) and full worldline flips (system-spanning
clusters when $\Jp$ is large), and dramatically improve mixing in stiff,
strongly-coupled regimes. \code{cluster\_sweeps} (default 0) counts these;
setting it $>0$ is referred to as the \emph{+SW} variant of an engine.

\begin{notebox}[Which moves to enable]
A robust general-purpose mix is
\code{sweeps\_per\_beta=20--60}, \code{worldline\_sweeps=3--5},
\code{cluster\_sweeps=0}; add \code{cluster\_sweeps=1} for strongly
frustrated instances or whenever convergence visibly stalls at late steps.
All three counters multiply runtime linearly.
\end{notebox}

\section{The annealing loop and replicas}

Putting it together, \code{SQAAnnealer.run()} executes, for each schedule
step $(\beta_k,\Gamma_k)$:
\begin{enumerate}
\item recompute $\Jp(\beta_k,\Gamma_k,M)$ (Eq.~\eqref{eq:jperp}) and the
slice temperature factor $\beta_k/M$;
\item run the configured number of slice, worldline, and cluster sweeps with
the acceptance rules of Section~\ref{sec:sqa-moves};
\item measure the energy of \emph{every slice} of every replica; record the
best classical projection seen so far.
\end{enumerate}

The \code{replicas} parameter runs $R$ independent copies of the whole
Trotter system inside one call (OpenMP-parallel, one RNG stream per replica).
Besides restart statistics, replicas are what powers the adaptive schedule's
susceptibility estimator (Chapter~\ref{ch:optimal}), which pools samples
across them.

At the end, the returned \code{best\_state} is the best single \emph{slice}
ever observed --- the ``classical projection'' of the quantum state --- with
\code{best\_energy} its true classical energy $\Ecl$, and
\code{energy\_trace} the per-step mean energy.

\section{Usage}

\subsection{One-call interface}

\begin{pycode}
from qanneal import solve

result = solve(
    ising,
    method="sqa",
    reads=16,             # independent restarts
    trotter_slices=32,    # M
    replicas=4,           # parallel Trotter systems per read
    sweeps_per_beta=60,
    worldline_sweeps=4,
    cluster_sweeps=1,     # enable SQA+SW
    seed=42,
)
print(result.best_energy, result.best_sample)
\end{pycode}

\subsection{Explicit schedules}

If no schedule is given, \code{solve()} builds a tuned one
(Section~\ref{sec:sqa-tuned}). To control it explicitly:

\begin{pycode}
import numpy as np
from qanneal import SQASchedule, SQAAnnealer

schedule = SQASchedule.from_vectors(
    betas  = np.linspace(0.1, 5.0, 80).tolist(),
    gammas = np.geomspace(5.0, 0.01, 80).tolist(),   # geometric decay!
)
ann = SQAAnnealer(ising, schedule, trotter_slices=32, replicas=4)
ann.set_seed(7)
res = ann.run(sweeps_per_beta=60, worldline_sweeps=4, cluster_sweeps=1)
print(res.best_energy)
\end{pycode}

\begin{warnbox}[Use geometric $\Gamma$ decay]
$\Jp$ depends on $\Gamma$ logarithmically at small $\Gamma$
(Eq.~\eqref{eq:jperp}), so a \emph{linear} $\Gamma$ ramp spends most of its
steps in a narrow band of $\Jp$ and rushes the physically important region.
Geometric decay (\code{np.geomspace}) distributes steps roughly uniformly in
$\ln\Gamma$, i.e.\ far more evenly in $\Jp$.
\end{warnbox}

\subsection{Problem-adaptive tuned SQA schedules}
\label{sec:sqa-tuned}

\code{auto\_schedule\_sqa\_tuned(problem, mode=...)} extends the SA tuner of
Section~\ref{sec:tuned-schedules}: the same random-flip probe estimates the
energy scale $\Delta$, then
\begin{equation}
\Gamma_{\text{start}}=m_\Gamma\,\Delta,\qquad
\Gamma_{\text{end}}=\Gamma_{\text{start}}\cdot r_\Gamma\cdot
\code{gamma\_end\_scale},
\end{equation}
with $m_\Gamma\in\{2.0,\,3.0,\,4.5\}$ for
\code{fast}/\code{balanced}/\code{accurate} and geometric spacing in between.
$\beta$ ramps exactly as in the SA tuner. Setting
$\Gamma_{\text{start}}$ a few times the flip scale guarantees the run
actually begins in the quantum-disordered phase.

\begin{pycode}
from qanneal import auto_schedule_sqa_tuned
schedule = auto_schedule_sqa_tuned(ising, mode="balanced")
\end{pycode}

\section{Parameter summary}

\begin{center}
\small
\begin{tabular}{@{}llp{8.6cm}@{}}
\toprule
\textbf{Parameter} & \textbf{Default} & \textbf{Meaning}\\
\midrule
\code{trotter\_slices} & 32 & $M$; more slices $\Rightarrow$ less Trotter error, linearly more memory/time.\\
\code{replicas} & 1 & Independent Trotter systems per run (OpenMP-parallel). Use $\geq4$ with the optimal schedule.\\
\code{sweeps\_per\_beta} & 20 & Single-spin sweeps per schedule step.\\
\code{worldline\_sweeps} & 5 & Whole-worldline flips per step; essential at large $\Jp$.\\
\code{cluster\_sweeps} & 0 & Swendsen--Wang segment flips per step ($>0$ = SQA+SW).\\
\code{continuous\_time\_slices} & 0 & If $>0$, overrides $M$ to approximate continuous time.\\
\bottomrule
\end{tabular}
\end{center}
